\subsection{Local projectives, Mittag-Leffler indices and tensor products}\label{algebra-proof-012-local-projectives-mittag-leffler-indices-and-tensor-products}

Separate editorial evidence for the Stacks Project authors' original \texttt{algebra.tex}, authority commit \texttt{a044\allowbreak{}46e5\allowbreak{}7ec1\allowbreak{}fbc2\allowbreak{}52a8\allowbreak{}71af\allowbreak{}cec7\allowbreak{}752f\allowbreak{}b280\allowbreak{}7b14}, lines 20589--21441. Authority SHA-256: \texttt{FA8B\allowbreak{}B92E\allowbreak{}58A4\allowbreak{}F78A\allowbreak{}2BD0\allowbreak{}1B3B\allowbreak{}6A4A\allowbreak{}87DE\allowbreak{}0A0D\allowbreak{}279F\allowbreak{}5DD9\allowbreak{}0641\allowbreak{}B574\allowbreak{}DD5F\allowbreak{}BFFF\allowbreak{}A4F3}. All units in this interval and its five received reports were read. The current interval differs by exactly the two already composed operations of \texttt{MC-STK-ERR-0440}; no other change or FAC environment occurs here. The following proposition, beginning at 21442, remains outside this adjudication.

Four minimal operations are proposed: two stage-choice repairs, one adverb, and one coefficient projection. The existing tensor correction retains its identity. Complete supporting arguments and the two further consequences below are editorial material, with the source's statements, hypotheses and exposition kept identifiable. No chapter or translation was mutated; a replay is not an instruction to replace a translation by these notes. No novelty or exhaustive literature claim is made.

\subsubsection{The original local-projective proof}\label{algebra-proof-012-the-original-local-projective-proof}

Source 20600--20693 receives the already reviewed decomposition into countably generated projective summands. Its proof is correct. Here are the exact matrix and summand checks; they generate no source operation.

Retain the original free module \(F=P\oplus Q\), basis, and shortest expression \(x=\sum_{i=1}^n a_i e_i\). Such a shortest expression exists because its possible finite lengths are natural numbers. If \(a_j=\sum_{i\ne j}a_i c_i\), the basis change \(e_i\mapsto e_i+c_i e_j\) for \(i\ne j\), fixing \(e_j\) and the remaining basis, has inverse obtained by subtracting these same multiples. In that basis the coefficient at \(e_j\) is zero. Thus no \(a_j\) belongs to the ideal generated by the other coefficients.

Write the original projections as \(e_i=y_i+z_i\), and \(y_i=\sum_j b_{ij}e_j+t_i\). The equality \(x=\sum_i a_i y_i\) gives, for each \(j\), \[
0=\sum_{i=1}^n a_i(b_{ij}-\delta_{ij}).
\] If any coefficient \(b_{ij}-\delta_{ij}\) in this relation were a unit, it would express the corresponding \(a_i\) in terms of the others. Every such coefficient lies in the maximal ideal \(\mathfrak m\). Consequently \(B=(b_{ij})\) reduces to the identity modulo \(\mathfrak m\), and \(\det B\) is a unit. The adjugate formula gives its actual inverse over \(R\).

In the original decomposition \(F=R^n\oplus F_{\rm rest}\), let \(T:R^n\to F_{\rm rest}\) have columns \(t_i\). The endomorphism sending \(e_i\) to \(y_i\) has the formula \[
(v,w)\longmapsto(B^{\mathsf t}v,Tv+w),
\qquad
(v,w)\longmapsto((B^{\mathsf t})^{-1}v,w-T(B^{\mathsf t})^{-1}v)
\] for its inverse. The transpose is required by the source's indexing of \(b_{ij}\). It follows that \(N=\sum_i Ry_i\) is a finite free summand of \(F\). Its retraction restricts to a retraction \(P\to N\), and \(x=\sum_i a_i y_i\in N\). If \(x=0\), the zero free summand suffices.

For a countably generated projective \(P\), every complement of a finite free summand is projective: restrict the projective lifting property using that complement's inclusion and retraction. At the next generator, the preceding construction gives a free summand of the complement containing its projected value. That value has finite support in a basis of this free summand; the span of those finitely many basis elements is again a direct summand, so it can be used as the next finite free piece. The partial sums are split and exhaust the chosen generators. The map from their external direct sum to \(P\) is injective because every finite partial sum is direct, and surjective because every generator occurs. Finally, the earlier countable-projective decomposition and the union of the resulting bases prove the original unrestricted local-ring theorem.

\subsubsection{Stable images and countable cofinality}\label{algebra-proof-012-stable-images-and-countable-cofinality}

Source 20697--20877. The following is a separate weakening of the index hypothesis: \textbf{a countable cofinal subset of the original directed set suffices} for both nonemptiness of the inverse limit of nonempty ML sets and exactness of the inverse limit of the stated short exact sequences. The source's countable-index statements are preserved.

For the original maps \(\varphi_{ji}:A_j\to A_i\), write \[
A'_i=\bigcap_{j\geq i}\varphi_{ji}(A_j).
\] If the images stabilize beyond \(j_0\geq i\), their stable value equals this intersection: for any other \(j\geq i\), choose \(k\geq j,j_0\). The stable value is the image from \(k\), hence is contained in the image from \(j\). This also handles incomparable indices in the definition of the intersection.

If \(j\geq i\) and \(a\in A'_j\), then for every \(k\geq i\) choose \(\ell\geq k,j\) and a preimage of \(a\) in \(A_\ell\). Its image in \(A_i\) factors through \(A_k\); thus \(\varphi_{ji}(a)\in A'_i\). These are the original restricted maps. They are surjective under the ML condition: choose \(k\geq j\) beyond stabilization thresholds for both \(i\) and \(j\). Any \(a_i\in A'_i\) has a preimage \(a_k\in A_k\); then \(\varphi_{kj}(a_k)\in A'_j\) maps to \(a_i\). If every \(A_i\) is nonempty, each stable image is nonempty. A compatible original family belongs coordinatewise to every image defining \(A'_i\), so inclusion induces an actual bijection \(\lim A'_i\to\lim A_i\), whose inverse retains the same coordinates.

Let \(C\subset I\) be a countable cofinal subset of a nonempty directed set. Enumerate it as \(c_1,c_2,\ldots\), repeating elements when finite, and choose \(j_1=c_1\), \(j_n\geq j_{n-1},c_n\). The sequence is cofinal in the original \(I\). Restriction of a compatible family gives the map to the limit over this sequence. Conversely, for a compatible sequence \((a_n\in A_{j_n})\), define \[
a_i=\varphi_{j_n i}(a_n)\quad\text{for any }n\text{ with }j_n\geq i.
\] Any two choices give the same value after passing to their larger sequence index. The resulting family is compatible: for \(i\leq k\), use a single \(j_n\geq k\) and the original composition rule. Both composites with restriction are identities. The same formulas are group or module homomorphisms in those categories and commute with morphisms of systems.

Apply these maps to the stable-image system. Starting with any element of \(A'_{j_1}\), successively choose a lift along each surjective map \(A'_{j_{n+1}}\to A'_{j_n}\), using the same axiom of choice as the source. The resulting compatible sequence gives a family over \(I\). This proves nonemptiness. If an empty index is allowed under a different convention, its set limit is a singleton and its group limit is zero, so the corresponding assertions hold separately without an enumeration.

For exactness retain the original short exact sequences and all their maps: \[
0\longrightarrow A_i\xrightarrow{f_i}B_i\xrightarrow{g_i}C_i\longrightarrow0.
\] Write \(\varphi^A,\varphi^B,\varphi^C\) for the three original transition families. Coordinatewise injectivity of \(f_i\) gives injectivity on limits. If a compatible \((b_i)\) maps to zero, its unique preimages \(a_i\) under \(f_i\) are compatible by commutation and injectivity, proving exactness in the middle.

Given \((c_i)\in\lim C_i\), put \(E_i=g_i^{-1}(c_i)\). Each is nonempty and the original \(B\)-maps restrict to them. Fix \(i\), and let \(j\) be a stabilization index for the images from \(A\). For \(k\geq j\), choose \(e_j\in E_j\), \(e'_k\in E_k\), and set \(e'_j=\varphi^B_{kj}(e'_k)\). There is a unique \(a_j\in A_j\) with \(f_j(a_j)=e_j-e'_j\). Stabilization gives \(a_k\in A_k\) with \(\varphi^A_{ki}(a_k)=\varphi^A_{ji}(a_j)\). Then \[
e_k=e'_k+f_k(a_k)\in E_k,
\qquad
\varphi^B_{ki}(e_k)=\varphi^B_{ji}(e_j).
\] The last equality follows by inserting the defining difference and the commutative square for \(f\). The reverse image inclusion follows by factorization through \(j\), so \(E\) is ML. The just-proved nonemptiness result gives \(\lim E_i\ne\varnothing\), proving surjectivity onto \((c_i)\). This proves exactness with a countable cofinal subset. For \(I=\mathbb N\), the criterion at 20865--20867 is exactly stabilization of these same images, so the source sequential lemma is an unchanged receiving case. No claim removes countable cofinality altogether.

\subsubsection{Domination and the original pushout}\label{algebra-proof-012-domination-and-the-original-pushout}

Source 20905--20996. The finite-test lemma is correct: for an arbitrary tensor-test module, express it as a filtered colimit of finitely presented modules. Any kernel element is represented at a stage and its image vanishes at a later stage. Apply the assumed kernel inclusion there and pass to the colimit. Conversely, the all-module condition includes those tests. This element argument preserves the original tensor maps and does not require a tensor functor to preserve arbitrary injections.

The original pushout is \[
N'=(M'\oplus N)/\{(g(x),-f(x)):x\in M\},
\quad f'(m')=[m',0],\quad g'(n)=[0,n].
\] The map \(\ker f\to\ker f'\) is \(x\mapsto g(x)\). It is onto because \([m',0]=0\) means \((m',0)=(g(x),-f(x))\) for some \(x\in\ker f\); its kernel is \(\ker f\cap\ker g\). After tensoring with any test module, the resulting pushout has the same quotient presentation for the tensor maps, by right exactness. Applying this exact kernel argument to that new pushout gives the source sequence. It is not obtained by declaring the old kernel sequence preserved under tensoring. This proves the domination/universal-injectivity equivalence.

The cokernel map \([m',n]\mapsto[n]\) identifies \(\operatorname{coker}f'\) with the original \(\operatorname{coker}f\). Under finite presentation, a universally injective \(f'\) has a retraction by the preceding source split criterion. If that retraction is \(h'\), then \(h=h'g'\) satisfies \(hf=g\). Conversely, if \(hf=g\), the formula \[
h'([m',n])=m'+h(n)
\] is well-defined since it kills \((g(x),-f(x))\), and is a retraction of \(f'\). This supplies the actual two directions needed below.

\subsubsection{Equalize the lifted maps at a finite stage}\label{algebra-proof-012-equalize-the-lifted-maps-at-a-finite-stage}

Source \texttt{proposition-ML-characterization}, 20999--21103; the missing step is at 21031. Keep its directed system \((M_i,f_{ij})\), colimit \(M\), and canonical maps \(f_i\). Condition (1) supplies \(g:M_i\to Q\) with exactly the same tensor kernels as \(f_i\). The cokernel of \(g\) is finitely presented: choose a finite presentation of \(Q\) and add lifts of the images of finitely many generators of \(M_i\) as finitely many extra relations. The domination lemma therefore gives \(h:Q\to M\) with \(hg=f_i\). Since \(Q\) is finitely presented, choose \(j\geq i\) and \(h':Q\to M_j\) with \(f_jh'=h\).

At this point only \[
f_j f_{ij}=f_i=f_jh'g
\] has been proved. Equality after the colimit map does not assert \(f_{ij}=h'g\) in \(M_j\). Choose generators \(a_1,\ldots,a_m\) of the original \(M_i\). Each difference \(f_{ij}(a_r)-h'g(a_r)\) maps to zero in the colimit, so vanishes after some transition from \(M_j\). Choose a common upper stage \(j'\) for these finitely many stages. Then \[
f_{ij'}=f_{jj'}h'g.
\] When the generating list is empty, both maps from \(M_i=0\) agree already. Replacing the stage by \(j'\) and the lift by \(f_{jj'}h'\) makes precisely the printed diagram commute. Its factorization shows that the new \(f_{ij'}\) dominates \(g\), which dominates \(f_i\), proving (2). The proposed insertion states this stage enlargement and the resulting equality, with the proof above retained separately.

For the converse, factor \(f:P\to M\), with \(P\) finitely presented, as \(f_i g'\). Choose \(j\) from (2). For every test module \(T\), \(f_i\) factors through \(f_{ij}\), so their tensor kernels on \(M_i\otimes T\) coincide by (2). Taking their inverse images under \(g'\otimes1\) gives exactly the equality of kernels for \(f\) and \(f_{ij}g'\). The latter map has finitely presented target \(M_j\), proving (1).

Conditions (2) and (3) are equivalent with the original quantifiers. If (2) holds, \(f_i\) dominates each \(f_{ik}\); transitivity and finite presentation of \(\operatorname{coker}f_{ik}\) give the required factorization of \(f_{ij}\) through \(f_{ik}\). Conversely, every element of \(\ker(f_i\otimes1_T)\) vanishes at some stage \(k\geq i\), by the original tensor-colimit map. A factorization through that stage makes it vanish under \(f_{ij}\otimes1_T\). This proves (2) and retains every tensor kernel rather than replacing them by a different presentation.

\subsubsection{Use a common upper bound for incomparable indices}\label{algebra-proof-012-use-a-common-upper-bound-for-incomparable-indices}

The second missing choice is at 21101--21102 in (5) implies (3). For (3) implies (4), fix \(i\) and the one index \(j\) from (3). If \(k\geq j\), a map \(M_j\to N\) composed with the factor \(M_k\to M_j\) gives the inclusion of its image in the image from \(M_k\). The reverse inclusion follows from \(f_{ik}=f_{jk}f_{ij}\). Thus the same index stabilizes the original Hom images for every \(N\). Condition (4) includes (5).

For (5), keep \(N=\prod_{s\in I}M_s\) and its stabilization index \(j\). There is a map \(e_j:M_j\to N\) equal to the identity in coordinate \(j\) and zero in all other coordinates. For \(k\geq j\), equality of the two Hom images gives \(H_k:M_k\to N\) such that \(H_k f_{ik}=e_j f_{ij}\). Projection to coordinate \(j\) gives \(h_k:M_k\to M_j\) with \[
f_{ij}=h_k f_{ik}.
\] This proves exactly the factor needed from the source's product-coordinate step. It does not require any assertion about exchanging arbitrary images with a product.

The printed final sentence treats \(k\geq j\) and \(k\leq j\); these need not exhaust a directed set. For any original \(k\geq i\), choose \(\ell\geq k,j\). The factor just proved for \(\ell\) gives \(h_\ell:M_\ell\to M_j\). Set \[
h=h_\ell f_{k\ell}:M_k\to M_j.
\] Then \(h f_{ik}=h_\ell f_{i\ell}=f_{ij}\). This proves (3) for incomparable indices as well. The proposition's five statements and definition of an ML module are unchanged.

\subsubsection{Retain the two tensor corrections and the finite-free test}\label{algebra-proof-012-retain-the-two-tensor-corrections-and-the-finite-free-test}

Source 21116--21274. Reports \texttt{OCC-12097}, \texttt{OCC-00452} and \texttt{OCC-00453} describe the two operations already composed as \texttt{MC-STK-ERR-0440}. The actual second system is \((N_j,g_{jj'})\). Thus the colimit terms are \(M_i\otimes_R N_j\), and the second factor map is \(b:N_{j''}\to N_{j'}\). With these types, \[
(a\otimes b)(f_{ii''}\otimes g_{jj''})
=f_{ii'}\otimes g_{jj'}.
\] The original tensor product is the colimit over the directed product set: its comparison sends the class of \(m_i\otimes n_j\) to \(f_i(m_i)\otimes g_j(n_j)\). For the inverse, choose representatives for both factors, take their tensor class, and use a common later pair to prove independence; additivity and balancing hold at a common pair. The two maps are inverse on simple tensors. The terms are finitely presented as in the source. Condition (3), now checked for every later pair including incomparable coordinates, proves the original tensor-closure result. No new operation duplicates either retained edit.

For the flat-module remark, write \(D_i=\operatorname{Hom}_R(M_i,R)\) for the duals of the original finite free stages, and let \(L_i\subset D_i\) be the stable image. Each sufficiently late map \(D_j\to D_i\) is the composite of a surjection onto \(L_i\) with its inclusion. After tensoring with an arbitrary \(T\), right exactness keeps that first map surjective; the resulting image is exactly \(\operatorname{im}(L_i\otimes T\to D_i\otimes T)\). This expression is independent of the later stage. The original finite-free Hom/tensor isomorphism carries it to the Hom images, proving the remark. Injectivity of \(L_i\otimes T\to D_i\otimes T\) is not needed.

Report \texttt{OCC-00454} only requires ``is clearly weaker'' at 21173. The converse argument is valid. If \(F\twoheadrightarrow P\) is the original finite free surjection and \(F\to Q\) has the same tensor kernels as \(F\to M\), its kernel contains that of \(F\to P\), so it induces \(g:P\to Q\). After tensoring, \(F\otimes T\to P\otimes T\) is onto. The inverse images under this surjection of the two kernels from \(P\otimes T\) are equal by hypothesis, hence the kernels themselves are equal. This proves the finite-free test without any additional source correction.

Restriction of scalars at 21200--21219 uses finitely presented \(S\)-stages that remain finitely presented over \(R\), under the stated finite and finitely presented ring-map hypothesis. The same actual factor maps remain \(R\)-linear. Although the paragraph only records factors for \(k\geq j\), the common-upper-bound construction above extends them to all \(k\geq i\), so no additional edit is necessary. In the quotient lemma 21221--21245, the same finite-presentation fact and the equality of the original \(R\)-linear and \(S\)-linear Hom sets justify the proof. A wider one-way implication is proved separately below.

The dual-number example keeps its original action \(\epsilon(a,b)=(0,u(a))\) on \(F_1\oplus F_0\); its square is zero. The quotient map \((a,b)\mapsto(a,[b])\) has kernel \(0\oplus\operatorname{im}u=\epsilon M\), so it induces exactly \(M/\epsilon M\cong F_1\oplus M_0\). Direct summands of ML modules are ML: given \(f:P\to L\), compose with a split inclusion \(L\to M\). Tensoring that inclusion remains injective because its retraction also tensors; hence the two maps from \(P\) have identical tensor kernels. Apply condition (1) for \(M\). Thus a non-ML summand \(M_0\) forces this quotient to be non-ML. Free modules are ML because any map from a finitely presented module has image in a finite coordinate summand, whose inclusion remains split after tensoring. If the original \(M\) were ML over \(S\), tensor closure with the finitely presented \(S/(\epsilon)\) would make its quotient ML over \(S\), and the quotient lemma would make it ML over \(R\), a contradiction. These arguments verify the original example without changing its action or presentation.

\subsubsection{The coefficient projection and the original product criteria}\label{algebra-proof-012-the-coefficient-projection-and-the-original-product-criteria}

Source 21277--21441; report \texttt{OCC-00455}. The proof at 21349 defines \(p_x:M^M\to M\), whereas \(a_i\) belongs to \(R^M\). Keep this \(p_x\) and replace only its out-of-domain application \(p_x(a_i)\) at 21353 by the coefficient \((a_i)_x\in R\). The canonical tensor-product map then gives exactly \[
x=p_x(d)=\sum_i p_x(f(x_i\otimes a_i))=\sum_i(a_i)_x x_i.
\] The finitely many original \(x_i\) generate \(M\). The zero module permits an empty sum. Conversely, a finite free surjection onto \(M\) tensors onto each receiving module and onto the product. Finite free tensoring commutes with the product by its coordinate maps, and the product of the surjections is onto by choice. The source square therefore proves the surjectivity criterion.

For finite presentation, keep a finite presentation \(R^m\xrightarrow{A}R^n\to M\to0\). The finite-free product comparison isomorphisms carry the image of the original matrix \(A\) in the first row of the source diagram to its image in the second row; those rows are right exact. Consequently they induce inverse isomorphisms on the two cokernels, precisely the canonical map for \(M\). In the converse, retain \(K=\ker(F\to M)\) for the chosen finite free surjection. For any set \(A\) and \(k\in K^A\), include it in \(F^A\), and use the isomorphism \(F\otimes R^A\to F^A\) to lift it to \(z\). The image of \(z\) in \(M\otimes R^A\) maps to zero in \(M^A\); injectivity of that comparison makes this image zero. Right exactness lifts \(z\) from \(K\otimes R^A\). By the commutative square and injectivity \(K^A\to F^A\), that lift maps to the original \(k\). Thus the comparison for \(K\) is onto for every \(A\), and the preceding criterion makes \(K\) finite. No injection at the left of the tensor row was assumed.

The original rational examples are correct. The residue map \(\mathbb Z\to\prod_{n\geq1}\mathbb Z/n\) is injective since an integer divisible by every positive integer is zero, including the zero factor at \(n=1\). Tensoring by the original localization \(\mathbb Q\) preserves that injection; each \(\mathbb Q\otimes\mathbb Z/n\) is zero. For the second example, every finite sum in \(\mathbb Q\otimes\prod\mathbb Z\) admits one common nonzero integer denominator, and conversely \((a_n/m)\) is the image of \((1/m)\otimes(a_n)\). The sequence \((1/n!)_{n\geq1}\) has no such denominator: it would require \(m/n!\) to be an integer for all \(n\), impossible for a fixed nonzero \(m\).

The example's non-ML premise can also be realized explicitly without using the as-yet-unreviewed tensor characterization. Write \(\mathbb Q\) as the colimit of \(M_n=\mathbb Z\), \(n\geq0\), with \(f_{n,n+1}\) multiplication by \(n+1\), and stage map \(m\mapsto m/n!\). The maps are compatible; every rational denominator divides some factorial, giving surjectivity to \(\mathbb Q\), and these injective stage maps give injectivity of the induced map from the colimit. For the test target \(\mathbb Z\), the image of \(\operatorname{Hom}(M_k,\mathbb Z)\) in \(\operatorname{Hom}(M_j,\mathbb Z)\) is \((k!/j!)\mathbb Z\). These subgroups do not stabilize: for each sufficiently large \(k\), the next is the proper multiple \((k+1)(k!/j!)\mathbb Z\). Condition (4) therefore proves \(\mathbb Q\) is non-ML. A presentation matching the original example is \(F_1\xrightarrow{u}F_0\to\mathbb Q\to0\), with bases \(c_n,b_n\), \[
u(c_n)=b_n-(n+1)b_{n+1},\qquad b_n\longmapsto1/n!.
\] The composites are zero. Any finite sum in \(F_0\) reduces modulo these relations to an integer multiple of the last basis element in its finite support, by successive substitutions. If its rational image is zero, that remaining integer is zero. This proves exactness and retains every coefficient of the presentation. It is an editorial realization of the source example, not a replacement of its general statement.

Finally, \texttt{lemma-kernel-tensored-fp}, 21414--21440, is correct. Factor the original map from the finitely presented \(P\) through some \(M_j\); the given single element \(x\) has image \(y\) in \(M_j\otimes Q\). Since its colimit image is zero, the defining filtered-colimit equivalence gives a later stage where that same \(y\) is zero. The original composite \(P\to M_j\to M_{j'}\) is the required map into the finitely presented \(P'=M_{j'}\). This uses a witness for the one element, not a claim that an entire arbitrary kernel vanishes at a uniform stage.

\subsubsection{Mittag-Leffler base change and descent along pure ring maps}\label{algebra-proof-012-mittag-leffler-base-change-and-descent-along-pure-ring-maps}

This is a separate consequence of the corrected characterization, receiving the finite-equation method of \texttt{ALGEBRA-RECON-501}. For every ring map \(R\to S\), an ML \(R\)-module \(M\) gives an ML \(S\)-module \(M\otimes_R S\). If the ring map is pure, meaning that every original unit map \[
X\longrightarrow X\otimes_R S,\quad x\longmapsto x\otimes1
\] is injective, the converse also holds. No flatness of \(S\) is assumed. Here is the complete factorization calculation.

First consider actual \(R\)-linear maps \(u:A\to B\), \(v:A\to C\), where \(A\) is finitely generated and \(B\) is finitely presented; \(C\) is arbitrary. Suppose there exists \(h_S:B\otimes S\to C\otimes S\) with \(h_Su_S=v_S\). Choose generators \(a_1,\ldots,a_m\) of \(A\) and a finite presentation \[
R^t\xrightarrow{D}R^n\xrightarrow{q}B\to0.
\] For each \(\ell\) choose an actual lift \(\sum_{r=1}^n U_{r\ell}e_r\) of \(u(a_\ell)\), retaining the entire matrices \(D=(D_{rs})\) and \(U=(U_{r\ell})\). Set \(c_\ell=v(a_\ell)\). Define \[
\begin{split}
\beta:C^n&\longrightarrow C^t\oplus C^m,\\
(z_r)_{r=1}^n&\longmapsto
\left(\left(\sum_{r=1}^n D_{rs}z_r\right)_{s=1}^t,
\left(\sum_{r=1}^n U_{r\ell}z_r\right)_{\ell=1}^m\right),\\
c&=(0,(c_\ell)_{\ell=1}^m).
\end{split}
\] The values \(z'_r=h_S(q(e_r)\otimes1)\) solve \(\beta_S(z')=c\otimes1\): the first equations are the original relations in \(B\), and the second are the factorization on the chosen generators of \(A\). Tensoring commutes with these finite direct sums and with the cokernel \(W=\operatorname{coker}\beta\). Hence the class of \(c\) maps to zero in \(W\otimes S\). Purity of its unit map makes that class zero already in \(W\). Choose \(z\in C^n\) with \(\beta(z)=c\). The first equations make \(e_r\mapsto z_r\) descend to an actual map \(h:B\to C\), and the second give \(hu=v\) on every generator of \(A\), hence on all of \(A\). Empty presentations and zero modules use the same maps with their zero dimensions. No identification of a Hom module after base change is required.

Now choose any original directed presentation \(M=\operatorname{colim}_i M_i\) by finitely presented \(R\)-modules. Right exactness gives the base-changed finite presentations and tensor-colimit comparison, so \(M\otimes S=\operatorname{colim}_i(M_i\otimes S)\) with its actual transition maps. If \(M\otimes S\) is ML, then for each \(i\) one index \(j\geq i\) satisfies condition (3) over \(S\): for every \(k\geq i\), \((f_{ij})_S\) factors through \((f_{ik})_S\). Apply the preceding calculation to \[
A=M_i,\quad B=M_k,\quad C=M_j,\quad u=f_{ik},\quad v=f_{ij}.
\] It supplies an \(R\)-linear factor for every such \(k\), keeping the same \(i,j\). Thus the original system satisfies (3) over \(R\), proving descent. Conversely, for an ML \(M\), tensor each actual factor in (3) with \(S\); the resulting factors prove the criterion for the base-changed system. This converse uses no purity assumption and proves arbitrary base-change ascent.

This calculation propagates to the source quotient lemma without changing its printed assumptions. For any ideal \(I\) and any \(S=R/I\)-module \(M\), the maps \[
M\otimes_R S\longrightarrow M,\quad m\otimes\overline r\longmapsto rm,
\qquad M\longrightarrow M\otimes_R S,\quad m\longmapsto m\otimes1
\] are well-defined \(S\)-linear inverses because \(IM=0\) and tensor balancing gives \(rm\otimes1=m\otimes\overline r\). Consequently the implication from ML over \(R\) to ML over \(R/I\) does not need finite generation of \(I\). The other direction in the source retains its finite-generation hypothesis and its already checked restriction proof.

There is no conflict with the dual-number counterexample: it concerns the original \(S\)-module with its specified square-zero action, whereas the base-change theorem concerns \(M_R\otimes_R S\). For the example with underlying \(M_R=F_1\oplus F_0\), the latter is free over \(S\) with the base-changed basis; the former has the explicit non-ML quotient just computed. The actual \(S\)-linear comparison is \[
\mu:M_R\otimes_R S\longrightarrow M,\quad m\otimes s\longmapsto s\cdot m.
\] It is onto, has an \(R\)-linear section \(m\mapsto m\otimes1\), and satisfies \(\mu(m\otimes\epsilon)=(0,u(a))\) for \(m=(a,b)\). Thus the precise relationship is a quotient of the base-changed free module; ML is not asserted to survive arbitrary quotients. This both supplies the connecting map and checks the source example against the broader result.

Pure maps include genuinely nonflat examples. For \(\mathbb Z\to\mathbb Z\times\mathbb Z/2\), \(n\mapsto(n,\overline n)\), first-coordinate projection \(r\) is a ring retraction. The map \(X\otimes S\to X\), \(x\otimes s\mapsto r(s)x\), retracts every unit map, proving purity. Multiplication by two kills the nonzero element \((0,\overline1)\) of \(S\), so tensoring the injection \(\mathbb Z\xrightarrow{2}\mathbb Z\) with \(S\) fails to preserve injectivity. This proves nonflatness. The ML descent statement therefore uses the same additional class of ring maps as group 501, with its full finite-equation comparison proved above. It asserts neither arbitrary projectivity nor effective descent of arbitrary module objects.

\subsubsection{Propagation and reading limits}\label{algebra-proof-012-propagation-and-reading-limits}

Groups 505 and 506 repair the two stage choices in the original characterization. Group 507 retains \texttt{MC-STK-ERR-0440} for all three reports and propagates the corrected all-indices quantifier into its tensor factors. Groups 508 and 509 handle the adverb and scalar coefficient. Group 510 records countable-cofinal nonemptiness and exactness with zero source operations. Group 511 records arbitrary base-change ascent and pure descent, receives group 501 and the corrected characterization, and checks both the one-way quotient consequence and the actual comparison to the source's dual-number module. Later uses remain to be reviewed in source order; the broader combined-results paper remains after core completion.

The new bounded disk query \texttt{Mittag-Leffler\ pure} returns a TeX-routing record for \emph{Higher-Rank Mathieu Opers, Toda Chain, and Analytic Langlands Correspondence}, a research PDF-routing record for \emph{Operators, Inequalities and Approximation}, and the local unpublished PDF-routing record \texttt{collatz\_path\_algebra}. All three remain unread and unused as proof evidence; no PDF fallback was performed. The initial display omitted the local layer by using a mismatched layer name; inspection of the saved raw result identified that third hit before adjudication. The topic snapshot preserves all three actual results together with earlier query history. The original Stacks TeX and the fully written comparisons above supply this batch's evidence; the lookup is not a novelty survey. Partial lookahead 21442--21480 remains unadjudicated.
